\section{Introduction}
\label{sec:intro}

Activation steering lets us control what a language model says without changing its weights. The idea is simple: compute a direction in activation space that corresponds to a concept (say, positive sentiment), and add a scaled version of that direction during inference. This approach, popularized by \citet{turner2023activation} and \citet{zou2023representation}, works well and requires no training. But where in the network should we intervene?

Nearly all steering work targets the \emph{residual stream}---the $\dmodel$-dimensional representation that flows through the transformer, accumulating contributions from attention heads and MLP sublayers. This is a natural choice: the residual stream is the model's ``main highway,'' and a single perturbation there affects all downstream computation.

But the residual stream is not the only option. Inside each MLP sublayer, activations pass through an intermediate space that is $4\times$ wider ($\dmlp = 4 \cdot \dmodel$). After the up-projection and nonlinearity, the model operates in a 3072-dimensional space (for GPT-2 Small) before projecting back down to 768 dimensions. This expanded space is where MLP neurons act as feature detectors \citep{cunningham2023sparse}, and recent work shows it can be decomposed into interpretable, sparse features \citep{dunefsky2024transcoders, shafran2025decomposing}. A natural question arises: \emph{is the MLP intermediate space a better target for activation steering?}

There are reasons to think it might be. The space is wider, so concepts may be less superimposed. Post-activation sparsity (from GELU/ReLU) could provide a more natural basis for intervention. And \citet{shafran2025decomposing} show that MLP-space features derived via semi-nonnegative matrix factorization outperform residual-stream baselines at causal steering---though they use learned features rather than simple mean-difference vectors.

We test this hypothesis directly. We compute mean-difference steering vectors \citep{turner2023activation} in three spaces within GPT-2 Small---the residual stream, the MLP intermediate (post-GELU) space, and the MLP output---and compare their effectiveness at steering sentiment. We evaluate steering efficiency (sentiment shift per unit of KL divergence), linear separability (probe accuracy), and the sparsity structure of each space.

Our main findings are:
\begin{itemize}[leftmargin=*,itemsep=0pt,topsep=0pt]
    \item We show that mean-difference steering in the MLP intermediate space is feasible but $2.6\times$ less efficient than residual stream steering, as measured by sentiment probability shift per unit of KL divergence.
    \item We find that the MLP intermediate space is \emph{less} linearly separable than the residual stream (76\% vs.\ 86\% probe accuracy), contradicting the intuition that a wider space should be more linear.
    \item We characterize the sparsity structure of MLP steering vectors: just 2.8\% of neurons (86 out of 3072) carry 50\% of the steering signal, compared to 43\% of dimensions needed in the residual stream.
    \item We demonstrate that this concentration makes MLP steering fragile---steered generations degrade to incoherence at lower intervention strengths than in the residual stream.
\end{itemize}

These results clarify when and why the MLP intermediate space is a useful intervention target. The residual stream is more efficient for general-purpose steering because it integrates information from all model components. The MLP space, by contrast, encodes specific features that concentrate the steering signal in few neurons---useful for targeted interventions, but fragile under the blunt instrument of mean-difference steering.
